% Generated draft: requires mathematical and source review.
\documentclass[11pt]{article}
\usepackage[margin=25mm]{geometry}
\usepackage{fontspec}
\setmainfont{Latin Modern Roman}
\usepackage{amsmath,amssymb,mathtools,unicode-math}
\setmathfont{Latin Modern Math}
\usepackage{enumitem,needspace,xurl,hyperref}
\setlength{\parindent}{0pt}
\setlength{\parskip}{5pt}
\setlength{\emergencystretch}{3em}
\newcommand{\N}{\mathbb{N}}
\newcommand{\Z}{\mathbb{Z}}
\newcommand{\Q}{\mathbb{Q}}
\newcommand{\R}{\mathbb{R}}
\newcommand{\C}{\mathbb{C}}
\newcommand{\F}{\mathbb{F}}
\newcommand{\A}{\mathbb{A}}
\newcommand{\PP}{\mathbb{P}}
\newcommand{\E}{\mathbb{E}}
\newcommand{\eps}{\varepsilon}
\newcommand{\dd}{\,\mathrm{d}}
\begin{document}
\section*{UnsolvedMath ID: 30005952}
\textbf{Source problem code:} catalogue entry
\subsection{Definitions and mathematical statement}
A rational homology three-sphere is a closed connected three-manifold \(Y\) satisfying \(H_*(Y;\Q)\cong H_*(S^3;\Q)\). The notation \(HF^-(Y)\) denotes minus Heegaard Floer homology, regarded as a module over \(\F[U]\), and \(HF^-_{\mathrm{red}}(Y)\) denotes its reduced part. Here \(U\) acts by the module structure. The supplied record does not specify the coefficient field \(\F\), orientation conventions, the precise reduced-homology convention, or whether direct summands must respect gradings.

For a nonzero reduced module annihilated by a power of \(U\), let
\[
\ell=\min\{n\geq 1: U^nHF^-_{\mathrm{red}}(Y)=0\}.
\]
For each positive integer \(j\), the module \(\F[U]/(U^j)\) is the quotient by the ideal generated by \(U^j\).

\textbf{Supplied question (strong geography restriction).} For every rational homology three-sphere \(Y\), does \(HF^-(Y)\) contain a direct summand
\[
\bigoplus_{j=1}^{\ell}\F[U]/(U^j)
=
\F[U]/(U^\ell)\oplus\F[U]/(U^{\ell-1})\oplus\cdots\oplus\F[U]/(U)?
\]
In the category of ungraded \(\F[U]\)-modules, this asks whether there exists a module \(M\) such that
\[
HF^-(Y)\cong\left(\bigoplus_{j=1}^{\ell}\F[U]/(U^j)\right)\oplus M.
\]
This expresses the ungraded reading only; a graded interpretation would require additional conventions concerning grading shifts.

\textbf{Unresolved qualifications.} The source defines \(\ell\) as the least integer annihilating the reduced part, without explicitly specifying the allowed integers. If \(HF^-_{\mathrm{red}}(Y)=0\), the least nonnegative exponent is \(0\), whereas the least positive exponent is \(1\); these give different interpretations of the requested summand. The record does not settle this boundary case, so the displayed positive-length formulation does not resolve the full universal question. The source also does not explicitly state the nilpotence and finiteness conventions underlying the definition of \(\ell\).
\subsection{Short English statement}
The supplied question asks whether the minus Heegaard Floer homology of every rational homology three-sphere contains cyclic \(\F[U]\)-summands of every length from \(1\) through the nilpotence length of its reduced part. Coefficients, grading conventions, and the case of vanishing reduced homology require clarification. The supplied open label and verification claims are recorded as curation, not independently established facts.
\subsection{Sources}
\begin{itemize}
\item Supplied local UnsolvedMath registry record, permanent ID \(30005952\), title ``Strong Geography for Rational Homology Spheres,'' source problem code \texttt{catalogue entry}, in the Oberwolfach Reports Open Problems collection.
\item Supplied report citation: \emph{Topologie (2025)}, Oberwolfach Reports, DOI \texttt{10.4171/owr/2024/34}, pp.\ 1913--1974. Supplied source URL and literature reference labelled ``original OWR report'': \url{https://doi.org/10.4171/owr/2024/34}. The record gives proposed year \(2024\) and an \texttt{owr-year:2024} tag, while the citation labels the report \(2025\); this discrepancy is unresolved.
\item The supplied literature assessment, dated exactly \texttt{2026-08-22}, labels the problem open and describes the record as a truncated table-of-contents fragment containing no identifiable mathematical claim to verify. This is supplied dated curation, not independently verified; its own qualification requests independent verification of time-sensitive status and cited claims.
\item The record separately marks the statement \texttt{corrected\_verified} and claims that it was checked against the official report and preserves its hypotheses, notation, and target. This supplied verification claim has not been independently checked here and differs from the dated literature assessment.
\item No proposer is identified. The original extracted text appends ``The stable chain rule for orthogonal calculus'' and ``Connor Malin (joint work with Niall Taggart)'' after the geography question. Their relationship to that question is unestablished, so those names are not assigned as its authors.
\end{itemize}
\end{document}
